\newpage
\ifdefined\useChapters
\chapter{O Peso do Silêncio}
\else
\chapter{}
\fi

Acordei antes do despertador, como se um barulho que não existia tivesse me tirado do sono. A luz que entrava pela janela era comum, mas havia algo estranho no ar — um vazio que não era bem dor, mas parecia uma ausência antiga, como quando sentimos saudade de um lugar onde nunca estivemos.

— Você vai perder a hora. Levanta — diz Lúcia da porta, com a voz meio rouca de quem também não dormiu direito, embora provavelmente não saiba o motivo.

Me arrasto até o banheiro. Tudo está igual, mas eu não estou. Sinto um peso atrás da nuca, uma sensação que insiste em me acompanhar desde que abri os olhos, como se tivesse deixado algo para trás — algo importante — e agora fosse tarde demais para buscar.

Saímos de casa juntos. Lúcia caminha rápido, como sempre, mas, de vez em quando, olha para mim com um estranhamento que tenta disfarçar. Talvez seja coisa minha. Ou talvez não.

No caminho para a escola, algumas pessoas me encaram mais do que deveriam. Não hostilmente, só... como se esperassem que eu dissesse algo. Ou fizesse algo. Mas ninguém parece saber exatamente o quê.

Eu também não.

\bigskip

No metrô, quando o trem balança, por um reflexo estendo a mão como se pudesse me segurar no ar. Como se o ar fosse tátil. É ridículo. E pior: sinto vergonha imediatamente, sem entender por quê.

Deve ser só cansaço.

Ou algo que minha mente insiste em esconder de mim.

\bigskip

Na escola, o corredor está barulhento, mas, por algum motivo, para mim tudo parece abafado, como se eu estivesse debaixo d'água. O professor que substitui Português está na porta, segurando uma prancheta.

— Você. — Ele aponta para mim. — Vai apresentar hoje. A turma precisa de exemplos. Não faça drama.

Automaticamente, começo a concordar. É o que sempre fiz. É mais fácil aceitar do que contrariar os outros, mais fácil agradar do que incomodar.

Mas, quando abro a boca, minha voz não sai.

Não por medo.

Por... desconforto.

Por algo que puxa de dentro, dizendo que eu não devo isso a ninguém.

Um calor sobe no peito, como se uma lembrança quisesse emergir mas estivesse trancada por dentro. Não vejo imagens. Não ouço vozes. Apenas sinto.

\emph{Você não precisa disso.}

Respirei.

E pela primeira vez na vida, disse:

— Não.

Não alto. Não agressivo. Não dramático.

Apenas um “não”.

Um que cabe inteiro em mim.

O professor abre a boca para responder, mas não sabe o que dizer. Ele esperava submissão. Recebeu limite.

Lúcia, mais atrás, me olha com um espanto orgulhoso, como se estivesse vendo um irmão que nunca conheceu.

Sigo andando. O chão parece menos pesado, como se meus passos, finalmente, fossem meus.

\bigskip

Depois das aulas, fico sentado no banco em frente à saída, observando as pessoas passarem. Tudo parece normal demais. Talvez seja isso que me incomoda. Talvez seja isso que me salva.

No meio da movimentação, vejo Larissa e Lúcia vindo pela calçada. Conversam baixo, rindo por algum motivo que não ouço. Elas não parecem amigas antigas; parecem pessoas que se encontraram num dia aleatório e, por alguma razão que nenhuma das duas saberia explicar, ficaram. 

Há um cuidado silencioso entre elas, como se as duas soubessem — mesmo sem lembrar — que já estiveram do mesmo lado de uma coisa grande, perigosa e impossível. Uma coisa da qual sobreviveram juntas.

Assim que chegam no ponto de ônibus, Larissa mexe na mochila e tira uma pequena caixa de madeira, envelhecida, com marcas de tempo e poeira acumulada nas bordas.

— Olha isso aqui — diz ela para Lúcia, mostrando o que encontrou.

Lúcia se inclina para ver. É uma fotografia antiga, quase sépia, de um homem jovem de uniforme militar. Um sorriso torto, um olhar cansado. Um rosto que não deveria dizer nada a nenhuma das duas.

Mas diz a Larissa.

— Quem é esse? — pergunta Lúcia.

— Não faço ideia — responde Larissa, com a sobrancelha franzida. — Mas quando achei essa foto hoje... senti uma coisa estranha aqui — ela toca o peito. — Como se eu tivesse raiva e saudade ao mesmo tempo. E nenhuma das duas faz sentido.

Lúcia observa a expressão dela, preocupada.

— Você quer guardar isso?

— Quero… não sei por quê. Mas quero. — Larissa fecha a caixinha com cuidado, como quem guarda algo que ainda não entende.

As duas sobem no ônibus juntas. Sentam no mesmo banco. E seguem.

A rua parece menos solitária quando elas passam.

\bigskip

O sol já está baixo quando volto caminhando para casa. A cidade parece respirar devagar, como se tentasse se lembrar de um ritmo mais simples. E, pela primeira vez em muito tempo, sinto que eu também posso respirar sem pedir permissão para o ar.

Mais tarde, deitado na cama, encaro o teto branco e vazio. Penso no meu dia. No “não” que finalmente saiu da minha boca. No jeito como o ar entrou mais leve depois disso.

Talvez o recomeço fosse isso.

Não um mundo novo.  
Não um apagamento perfeito.  
Não um salvador brilhando no céu.

Talvez fosse a coragem singela de não fugir de mim.

Talvez fosse, pela primeira vez, escolher o que eu quero carregar —  
e o que eu não quero mais deixar que carreguem por mim.